\section{Mathematical Formulation}
\label{sec:theory}

\subsection{Model Geometry and Assumptions}
\label{sec:geometry}

The human urinary bladder is idealised as a thin spherical shell of radius
$R$, wall thickness $h$, completely filled with an inviscid, irrotational
fluid of density $\rho_f$ and bulk modulus $K_f$.  Unlike the abdominal
cavity---which requires an oblate spheroidal geometry with aspect ratio
$c/a \approx 0.67$~\cite{Mace2026browntone}---the bladder is approximately
spherical at moderate-to-full fill states, making the spherical shell
approximation both simpler and more appropriate.

The thin-shell assumption requires $h/R \ll 1$.  For the bladder at
\SI{300}{\milli\litre}, $h \approx \SI{1.8}{\milli\metre}$ and
$R \approx \SI{42}{\milli\metre}$, giving $h/R \approx 0.04$, which
satisfies the condition comfortably.  At low fill volumes, however, the
ratio increases: $h/R \approx 0.12$ at \SI{100}{\milli\litre} and
$h/R \approx 0.22$ at \SI{50}{\milli\litre}, exceeding the conventional
$h/R < 0.1$ threshold.  Results below \SI{150}{\milli\litre} should
therefore be treated with caution; they are retained in
Table~\ref{tab:frequencies} for completeness but marked with a dagger.
The Love--Kirchhoff hypothesis (normals remain straight and normal to
the deformed mid-surface) is applicable for fill states above
\SI{150}{\milli\litre}.

The bladder wall is modelled as a Kelvin--Voigt viscoelastic material with
complex modulus
\begin{equation}\label{eq:Estar}
  E^* = E(1 + i\eta),
\end{equation}
where $\eta$ is the loss tangent.  The quality factor is $Q = 1/\eta$ and
the viscous damping ratio is $\zeta = \eta/2$.  Published values of $\eta$
for bladder wall tissue range from 0.3 to
0.5~\cite{VanMastrigt1991,Barnes2016}; we adopt $\eta = 0.4$ ($Q = 2.5$,
$\zeta = 0.2$) as a central estimate.  The flexural rigidity of the shell is
\begin{equation}\label{eq:D}
  D = \frac{Eh^3}{12(1 - \nu^2)}.
\end{equation}


\subsection{Fill-Dependent Parameters}
\label{sec:fill}

All geometric and material properties vary with fill volume $V$
(\si{\milli\litre}).  We adopt the following constitutive relationships,
drawn from the literature reviewed in \S\ref{sec:introduction}.

\paragraph{Radius.}
The equivalent sphere radius follows directly from the fill volume:
\begin{equation}\label{eq:radius}
  R = \left(\frac{3V}{4\pi}\right)^{1/3},
\end{equation}
giving $R = \SI{2.3}{\centi\metre}$ at \SI{50}{\milli\litre} and
$R = \SI{4.9}{\centi\metre}$ at \SI{500}{\milli\litre}.

\paragraph{Wall thickness.}
We assume approximately constant tissue volume $V_\mathrm{tissue}$,
calibrated to $h_0 = \SI{5}{\milli\metre}$ at $V_0 = \SI{50}{\milli\litre}$.
The tissue shell volume is
\begin{equation}\label{eq:tissue_vol}
  V_\mathrm{tissue} = \frac{4}{3}\pi\left[(R_0 + h_0)^3 - R_0^3\right],
\end{equation}
from which $h(V)$ is obtained by solving
\begin{equation}\label{eq:h_solve}
  (R + h)^3 = R^3 + \frac{3 V_\mathrm{tissue}}{4\pi},
\end{equation}
subject to a physiological floor of $h_\mathrm{min} = \SI{1.5}{\milli\metre}$.
This yields $h = \SI{5.0}{\milli\metre}$ at \SI{50}{\milli\litre}, thinning
to \SI{1.5}{\milli\metre} at \SI{400}{\milli\litre} and above.

\paragraph{Elastic modulus.}
Bladder wall stiffness increases dramatically with stretch (strain-stiffening),
a hallmark of collagenous soft tissue~\cite{Fung1993}.  We model this as a
log-linear interpolation between vibrometry-derived bounds.

Nenadic et al.~\cite{Nenadic2013} report \emph{shear} moduli
$\mu = \SIrange{9.6}{48.7}{\kilo\pascal}$ from ultrasound bladder vibrometry.
For nearly incompressible tissue ($\nu = 0.49$), the Young's modulus is
\begin{equation}\label{eq:E_mu}
  E = 2(1 + \nu)\,\mu \approx 2.98\,\mu,
\end{equation}
giving $E \approx \SIrange{29}{145}{\kilo\pascal}$.  These converted bounds
are consistent with the tensile-test data of Barnes~\cite{Barnes2016}, who
reports $E = \SIrange{10}{800}{\kilo\pascal}$ depending on tissue layer and
stretch ratio.  We adopt:
\begin{equation}\label{eq:E_fill}
  E(V) = E_\mathrm{min} \left(\frac{E_\mathrm{max}}{E_\mathrm{min}}\right)^{\phi},
  \qquad \phi = \frac{V - V_\mathrm{min}}{V_\mathrm{max} - V_\mathrm{min}},
\end{equation}
where $E_\mathrm{min} = \SI{29}{\kilo\pascal}$ at $V_\mathrm{min} =
\SI{50}{\milli\litre}$ and $E_\mathrm{max} = \SI{145}{\kilo\pascal}$ at
$V_\mathrm{max} = \SI{500}{\milli\litre}$, with $\phi$ clipped to $[0, 1]$.

\paragraph{Intravesical pressure.}
The pressure--volume relationship follows the characteristic cystometric
curve~\cite{Griffiths1980}: a compliant phase at low fill followed by a
steep rise.  We model this as a quadratic:
\begin{equation}\label{eq:P_fill}
  P(V) = P_\mathrm{min} + (P_\mathrm{max} - P_\mathrm{min})\,\phi^2,
\end{equation}
where $P_\mathrm{min} = \SI{490}{\pascal}$ ($\approx \SI{5}{cmH_2O}$) and
$P_\mathrm{max} = \SI{2940}{\pascal}$ ($\approx \SI{30}{cmH_2O}$).

\begin{table}[htbp]
  \centering
  \caption{Bladder model parameters at representative fill states.}
  \label{tab:params}
  \begin{tabular}{lcccccc}
    \toprule
    Parameter & Symbol & 50~mL & 150~mL & 300~mL & 500~mL & Unit \\
    \midrule
    Radius & $R$ & 2.29 & 3.30 & 4.15 & 4.92 & cm \\
    Wall thickness & $h$ & 5.00 & 2.73 & 1.79 & 1.50 & mm \\
    Elastic modulus & $E$ & 29.0 & 41.5 & 70.9 & 145.0 & kPa \\
    Intravesical pressure & $P$ & 5.0 & 6.2 & 12.7 & 30.0 & cmH$_2$O \\
    Poisson's ratio & $\nu$ & \multicolumn{4}{c}{0.49 (all fill states)} & --- \\
    Wall density & $\rho_w$ & \multicolumn{4}{c}{1050} & kg/m$^3$ \\
    Fluid density & $\rho_f$ & \multicolumn{4}{c}{1020} & kg/m$^3$ \\
    Fluid bulk modulus & $K_f$ & \multicolumn{4}{c}{2.2} & GPa \\
    Loss tangent & $\eta$ & \multicolumn{4}{c}{0.40} & --- \\
    \bottomrule
  \end{tabular}
\end{table}

Table~\ref{tab:params} summarises the model parameters at four representative
fill states.  The Poisson's ratio $\nu = 0.49$ reflects the near-incompressibility
of soft tissue.  The fluid properties correspond to urine at body temperature,
which is effectively identical to water for acoustic purposes.


\subsection{Breathing Mode ($n = 0$)}
\label{sec:breathing}

The breathing mode is a spatially uniform radial oscillation.  As in the
abdominal model~\cite{Mace2026browntone}, the restoring force is dominated
by the fluid bulk modulus:
\begin{equation}\label{eq:f0}
  f_0 = \frac{1}{2\pi}\sqrt{\frac{k_\mathrm{shell} + k_\mathrm{fluid}}
  {\rho_w h + \rho_f R}},
\end{equation}
where $k_\mathrm{shell} = 2Eh / [R^2(1-\nu)]$ and
$k_\mathrm{fluid} = 3K_f/R$.  For the bladder at \SI{300}{\milli\litre},
$k_\mathrm{fluid} \approx \SI{1.6e11}{\pascal\per\metre}$ while
$k_\mathrm{shell} \approx \SI{2.2e4}{\pascal\per\metre}$---a ratio
exceeding $10^6$.  The breathing mode frequency is approximately
\SI{9500}{\hertz}, firmly in the ultrasonic range and irrelevant to
WBV exposure.


\subsection{Flexural Modes ($n \geq 2$)}
\label{sec:flexural}

Flexural modes involve shape changes that preserve the enclosed volume to
leading order.  The fluid contributes only added mass, not volumetric
stiffness.  Following Lamb~\cite{lamb1882} and Junger \&
Feit~\cite{Junger1986}, the natural frequency of flexural mode $n$ is
\begin{equation}\label{eq:fn}
  \omega_n^2 = \frac{K_\mathrm{bend} + K_\mathrm{memb} + K_P}
                    {m_\mathrm{eff}^{(n)}},
  \qquad f_n = \frac{\omega_n}{2\pi},
\end{equation}
where the three stiffness contributions are:

\paragraph{Bending stiffness.}
\begin{equation}\label{eq:Kbend}
  K_\mathrm{bend} = \frac{n(n-1)(n+2)^2 D}{R^4}.
\end{equation}

\paragraph{Membrane stiffness.}
\begin{equation}\label{eq:Kmemb}
  K_\mathrm{memb} = \frac{Eh}{R^2}\,\lambda_n,
  \qquad \lambda_n = \frac{n^2 + n - 2 + 2\nu}{n^2 + n + 1 - \nu}.
\end{equation}

\paragraph{Pre-stress stiffness.}
The intravesical pressure $P$ produces a pre-stress in the shell membrane
that stiffens flexural modes:
\begin{equation}\label{eq:Kpre}
  K_P = \frac{P}{R}(n-1)(n+2).
\end{equation}

The effective mass per unit area includes both the shell wall and the fluid
added mass:
\begin{equation}\label{eq:meff}
  m_\mathrm{eff}^{(n)} = \rho_w h + \frac{\rho_f R}{n}.
\end{equation}

For a bladder at \SI{300}{\milli\litre}, the model predicts
$f_2 = \SI{13.9}{\hertz}$, $f_3 = \SI{24.1}{\hertz}$, and
$f_4 = \SI{35.3}{\hertz}$.  The $n = 1$ mode corresponds to rigid-body
translation and has zero natural frequency for a free shell; it is therefore
excluded from the analysis.


\subsection{Airborne Acoustic Coupling}
\label{sec:airborne}

At $f_2 \approx \SI{14}{\hertz}$ the acoustic wavelength in air is
$\lambda \approx \SI{25}{\metre}$, so the dimensionless wavenumber is
\begin{equation}\label{eq:ka}
  ka = \frac{2\pi f_2 R}{c_\mathrm{air}} \approx 0.011 \ll 1.
\end{equation}
The bladder lies deep in the Rayleigh regime.  In a multipole expansion
of the incident plane wave, the $n$-th order component scales
as~\cite{Junger1986}
\begin{equation}\label{eq:peff}
  p_\mathrm{eff}^{(n)} = p_\mathrm{inc} \times (ka)^n.
\end{equation}
For the $n = 2$ mode:
\begin{equation}\label{eq:peff2}
  p_\mathrm{eff}^{(2)} = p_\mathrm{inc} \times (ka)^2
  \approx p_\mathrm{inc} \times 1.1 \times 10^{-4}.
\end{equation}
Only 0.01\% of the incident pressure drives the fundamental flexural mode
via the airborne pathway.


\subsection{Mechanical (Whole-Body Vibration) Coupling}
\label{sec:wbv}

Whole-body vibration transmits through the skeleton and pelvic floor as base
excitation of the bladder wall.  The coupling chain comprises two stages:

\paragraph{Seat-to-pelvis transmissibility.}
ISO~2631 data indicate a pelvic resonance near
$f_\mathrm{pelvis} \approx \SI{5.5}{\hertz}$ with $\zeta_\mathrm{pelvis}
\approx 0.25$.  The transmissibility at frequency $f$ is
\begin{equation}\label{eq:T_pelvis}
  T(f) = \sqrt{\frac{1 + (2\zeta_p r_p)^2}{(1 - r_p^2)^2 + (2\zeta_p r_p)^2}},
  \qquad r_p = \frac{f}{f_\mathrm{pelvis}}.
\end{equation}
At $f_2 = \SI{13.9}{\hertz}$, $T \approx 0.29$: the pelvic system
attenuates the input, as we are above its resonance.

\paragraph{Bladder modal amplification.}
The bladder wall response to base excitation is characterised by the
single-degree-of-freedom frequency response:
\begin{equation}\label{eq:H_bladder}
  H(f) = \frac{1}{\sqrt{(1 - r^2)^2 + (2\zeta r)^2}},
  \qquad r = \frac{f}{f_n}.
\end{equation}
At resonance ($r = 1$), $H = Q = 2.5$.

\paragraph{Coupling ratio.}
The mechanical-to-airborne coupling ratio at the $n = 2$ resonance frequency is
\begin{equation}\label{eq:R_coupling}
  \mathcal{R}_\mathrm{bladder} = \frac{T(f_2) \times H(f_2)}{(ka)^2}
  \approx \frac{0.29 \times 2.5}{1.12 \times 10^{-4}}
  \approx 6.4 \times 10^3.
\end{equation}
Mechanical vibration transmitted through the pelvic skeleton is approximately
$3.8$ orders of magnitude more effective at exciting bladder wall modes than
airborne sound.
